% Kosaraju 两遍DFS。弧：1->2, 2->1, 2->3, 3->4, 4->3, 5->6, 6->4, 6->7, 7->1, 7->5。
% 第一遍（按编号重启、出邻居递增）：发现/完成 1:1/8, 2:2/7, 3:3/6, 4:4/5, 5:9/14, 6:10/13, 7:11/12；
% 完成序 4,3,2,1,7,6,5。第二遍在反向图按 5,6,7,1,2,3,4 启动：{5,7,6}、{1,2}、{3,4}。脚本见 scratchpad/ch2/kos2.py。
\begin{tikzpicture}[font=\figurefont\fontsize{8}{10}\selectfont,>=Stealth,
  vtx/.style={circle,draw=FigureStroke,line width=.8pt,minimum size=6.2mm,inner sep=0pt,text=FigureInk,fill=white},
  arc/.style={draw=FigureStroke,line width=.8pt,->,shorten >=.5pt},
  cross/.style={draw=FigureStroke,line width=1.1pt,->},
  tm/.style={font=\figurefont\fontsize{7.5}{9}\selectfont,text=FigureInk},
  grp/.style={draw=FigureMuted,line width=.6pt,dashed,rounded corners=1.8pt},
  note/.style={font=\FigureNoteFont\fontsize{7.5}{9}\selectfont,text=FigureMuted,align=center}]
\def\ctwoverts#1{%
  \node[vtx,#1] (v5) at (0,0) {5};
  \node[vtx,#1] (v7) at (0.95,0.95) {7};
  \node[vtx,#1] (v6) at (0.95,-0.95) {6};
  \node[vtx,#1] (v1) at (2.35,0.95) {1};
  \node[vtx,#1] (v2) at (3.55,0.95) {2};
  \node[vtx,#1] (v4) at (2.35,-0.95) {4};
  \node[vtx,#1] (v3) at (3.55,-0.95) {3};}
% ---------- (a) 原图与完成时刻 ----------
\begin{scope}
  \ctwoverts{}
  \draw[arc] (v5)--(v6); \draw[arc] (v6)--(v7); \draw[arc] (v7)--(v5);
  \draw[arc,bend left=18] (v1) to (v2); \draw[arc,bend left=18] (v2) to (v1);
  \draw[arc,bend left=18] (v3) to (v4); \draw[arc,bend left=18] (v4) to (v3);
  \draw[cross] (v7)--(v1); \draw[cross] (v2)--(v3); \draw[cross] (v6)--(v4);
  \node[tm,left=1pt of v5] {9/14};
  \node[tm,above=1pt of v7] {11/12};
  \node[tm,below=1pt of v6] {10/13};
  \node[tm,above=1pt of v1] {1/8};
  \node[tm,above=1pt of v2] {2/7};
  \node[tm,below=1pt of v4] {4/5};
  \node[tm,below=1pt of v3] {3/6};
  \node at (1.8,-2.3) {(a) 第一遍：原图上的发现/完成时刻};
\end{scope}
% ---------- (b) 反向图上的第二遍 ----------
\begin{scope}[xshift=5.15cm]
  \fill[FigureBlueFill,rounded corners=3pt] (-0.45,0) -- (0.95,1.42) -- (1.42,0.95) -- (1.42,-0.95) -- (0.95,-1.42) -- cycle;
  \fill[FigureGreenFill,rounded corners=3pt] (1.95,0.55) rectangle (3.95,1.4);
  \fill[FigureYellowFill,rounded corners=3pt] (1.95,-1.4) rectangle (3.95,-0.55);
  \ctwoverts{fill=white}
  \draw[arc] (v6)--(v5); \draw[arc] (v7)--(v6); \draw[arc] (v5)--(v7);
  \draw[arc,bend left=18] (v1) to (v2); \draw[arc,bend left=18] (v2) to (v1);
  \draw[arc,bend left=18] (v3) to (v4); \draw[arc,bend left=18] (v4) to (v3);
  \draw[cross] (v1)--(v7); \draw[cross] (v3)--(v2); \draw[cross] (v4)--(v6);
  \node[note] at (0.45,1.7) {第1次：根 5};
  \node[note] at (2.95,1.7) {第2次：根 1};
  \node[note] at (2.95,-1.7) {第3次：根 3};
  \node at (1.8,-2.3) {(b) 第二遍：反向图，按 5,6,7,1,2,3,4 启动};
\end{scope}
% ---------- (c) 缩点图 ----------
\begin{scope}[xshift=10.4cm]
  \node[vtx,minimum width=13mm,fill=FigureBlueFill,shape=rectangle,rounded corners=1.8pt,minimum height=6mm] (CS) at (0,0) {$\{5,6,7\}$};
  \node[vtx,shape=rectangle,rounded corners=1.8pt,minimum width=11mm,minimum height=6mm,fill=FigureGreenFill] (CM) at (1.9,0.95) {$\{1,2\}$};
  \node[vtx,shape=rectangle,rounded corners=1.8pt,minimum width=11mm,minimum height=6mm,fill=FigureYellowFill] (CT) at (1.9,-0.95) {$\{3,4\}$};
  \draw[cross] (CS)--(CM); \draw[cross] (CM)--(CT); \draw[cross] (CS)--(CT);
  \node[note,below=2pt of CS] {最大完成 14};
  \node[note,right=2pt of CM] {8};
  \node[note,right=2pt of CT] {6};
  \node at (1.0,-2.3) {(c) 缩点图};
\end{scope}
\end{tikzpicture}
